% Lab 7: Lab test. Compile: latexmk -xelatex lab07.tex
\documentclass[10pt,aspectratio=169,t]{beamer}
\usetheme[lab]{upbminimal}

\course{Mobile and Embedded Computing}
\decklabel{Lab 7}
\title{Lab test}
\subtitle{Rules, format, and two practice specs}
\author{Dinu-Ștefan Rusu}
\institute{FILS, Universitatea Politehnica București}
\date{Week 14}

\begin{document}

\titleframe

\begin{frame}[eyebrow=Today]{What this lab is for}
  \vskip 2mm
  \begin{grid}[2]
    \cell[\faIcon{stopwatch}]{Prove three skills, once}{Build a small app from a one-page spec in ninety minutes: a device message, an offline write, and a remote read.}
    \cell[\faIcon{graduation-cap}]{Worth 1.5 points}{0.5 points per requirement, graded from your pushed branch. There is no live demo.}
  \end{grid}
  \vskip 6mm
  \taskmeta{Time}{90 minutes}{Submit}{Push to your project repo, branch \code{lab-7}, before time is up}
\end{frame}

\begin{frame}[eyebrow=Rules]{Individual, timed, your own machine}
  \begin{steps}
    \item You work alone. No pair programming, no shared screens, no copying a teammate's branch.
    \item Ninety minutes, one sitting, starting when you receive the spec.
    \item Your own laptop, your existing project repository from Labs 1 to 6.
    \item A real board or the Wokwi simulator, whichever you used in the labs. Either is accepted.
  \end{steps}
  \begin{takeaway}{The clock does not pause.}
    A crashed simulator or a broker that refuses a connection eats into your ninety minutes. Restart calmly and keep going.
  \end{takeaway}
\end{frame}

\begin{frame}[eyebrow=Rules]{What you may and may not use}
  \begin{steps}
    \item Internet access is allowed: official docs, package pages, your own past labs.
    \item AI agents are allowed for writing code.
    \item You must be able to explain any line the grader points to, in your own words.
    \item Asking another person, in the room or online, to write code for you is not allowed.
  \end{steps}
  \begin{warning}
    An agent-written line you cannot explain earns no credit for that requirement, even if it compiles.
  \end{warning}
\end{frame}

\begin{frame}[embedded,eyebrow=Setup]{Infrastructure for the test}
  \vskip 2mm
  \begin{grid}[2]
    \cell[\faWifi]{A broker on the instructor's machine}{A Mosquitto broker for the MQTT requirement. No need to run your own.}
    \cell[\faServer]{Two Go servers on the instructor's machine}{The outbox server from Lab 5 and the gRPC server from Lab 6, already running.}
  \end{grid}
  \vskip 6mm
  \muted{Addresses and ports are announced at the start of the test. Write them down before the clock starts.}
\end{frame}

\begin{frame}[eyebrow=Format]{One spec, three requirements}
  \vskip 2mm
  \begin{grid}[2]
    \cell[\faIcon{file-alt}]{One page, handed at the start}{An app idea and three requirements, on paper or a shared document. Nothing changes once the clock starts.}
    \cell[\faClock]{Ninety minutes, one push}{Build, commit as you go, and push to branch lab-7 before time is up. No extensions.}
  \end{grid}
  \vskip 6mm
  \muted{Every spec asks for the same three kinds of requirement, described next.}
\end{frame}

\begin{frame}[embedded,eyebrow=Format]{The three requirements, generically}
  \begin{grid}[3]
    \cell[\faWifi]{Publish or subscribe over MQTT}{Talk to the broker: publish a reading, or subscribe and show a value live, with reconnect handled.}
    \cell[\faDatabase]{One offline-capable write}{A local write goes through an outbox and reaches the server once connectivity returns.}
    \cell[\faBluetoothB]{One gRPC or BLE read}{A unary call to the Go server, or a characteristic read from a peripheral, parsed into a model.}
  \end{grid}
\end{frame}

\begin{frame}[eyebrow=Rubric]{0.5 points per requirement}
  {\usebeamerfont{small}\begin{tabular}{@{}p{30mm}p{42mm}p{42mm}@{}}
    \toprule
    \thead{Requirement} & \thead{Full credit (0.5 pt)} & \thead{Partial credit (0.25 pt)} \\
    \midrule
    MQTT pub or sub & Connects, and reconnects after a broker drop & Connects once, but does not recover from a drop \\
    Offline write + outbox & Written offline, reaches the server once online & Outbox exists, but the drain loop never runs \\
    gRPC or BLE read & Parses the real response into a model, shown live & Call or read succeeds, but fields come from a raw map \\
    \bottomrule
  \end{tabular}}
  \vskip 2mm
  \muted{Anything not attempted scores zero.}
\end{frame}

\begin{frame}[fragile,eyebrow=Submission]{Push before the clock runs out}
\begin{shell}
git checkout -b lab-7
git push -u origin lab-7
\end{shell}
  \begin{steps}
    \item Create the branch the moment you receive the spec.
    \item Commit as you finish each requirement, with a message that names it.
    \item Push before the ninety minutes end.
    \item Open one pull request titled with your name. It does not need a review to count.
  \end{steps}
  \begin{warning}
    The grader reads the branch as it stands at the deadline. A late push is not graded.
  \end{warning}
\end{frame}

\begin{frame}[eyebrow=Prepare]{How to prepare}
  \vskip 2mm
  {\usebeamerfont{small}\begin{tabular}{@{}p{62mm}p{20mm}@{}}
    \toprule
    \thead{Requirement} & \thead{Taught in} \\
    \midrule
    Publish or subscribe over MQTT & Lab 2 \\
    One offline-capable write with an outbox & Lab 5 \\
    A gRPC unary call, or a BLE characteristic read & Lab 6, Lab 3 \\
    \bottomrule
  \end{tabular}}
  \vskip 4mm
  \begin{takeaway}{Nothing on the test is new.}
    Re-read the Task slides and the reference code in those three labs. Redo one small piece from memory before you come.
  \end{takeaway}
\end{frame}

\begin{frame}[embedded,eyebrow=Practice A]{Practice A: greenhouse monitor}
  \muted{An app that shows a greenhouse reading live and logs manual care events.}
  \vskip 2mm
  \begin{columns}[T]
    \begin{column}{0.56\textwidth}
      \begin{steps}
        \item Subscribe with \code{mqtt\_client} to \code{greenhouse/temp} and show the value live in a BlocBuilder, with reconnect on broker loss.
        \item A \emph{Log watering} button writes to a local table and an outbox in one transaction. A drain loop posts it to the outbox server.
        \item Call the gRPC server for the care thresholds for the plant and show them once parsed into a model.
      \end{steps}
    \end{column}
    \begin{column}{0.4\textwidth}
      \kv{Grader checks}{}
      \begin{checklist}
        \item Stopping and restarting the broker resumes live updates.
        \item A watering event logged in airplane mode reaches the server once online.
        \item The threshold values on screen come from the model, not a raw map.
      \end{checklist}
    \end{column}
  \end{columns}
\end{frame}

\begin{frame}[embedded,eyebrow=Practice B]{Practice B: delivery tracker}
  \muted{An app that tracks a van's status and one delivery scan.}
  \vskip 2mm
  \begin{columns}[T]
    \begin{column}{0.56\textwidth}
      \begin{steps}
        \item Publish the van's status to \code{fleet/van1/status} every few seconds with QoS 1.
        \item Marking a delivery complete writes to a local table and an outbox in one transaction, synced by a drain loop.
        \item Connect to the scanner peripheral over BLE and subscribe to the last scanned code, shown live on screen.
      \end{steps}
    \end{column}
    \begin{column}{0.4\textwidth}
      \kv{Grader checks}{}
      \begin{checklist}
        \item Publishing resumes after the broker connection drops and comes back.
        \item A delivery marked complete offline appears on the server once the sync loop runs.
        \item The scanned code updates the screen without polling the peripheral.
      \end{checklist}
    \end{column}
  \end{columns}
\end{frame}

\begin{frame}[eyebrow=Troubleshooting]{Mistakes that cost time}
  \vskip 2mm
  \trouble{Client not connected}{You published before the MQTT connect finished. Await the connection state before the first publish or subscribe.}
  \trouble{An outbox row stays queued forever}{The drain loop never marks the row sent after a successful post. Update its status inside the same transaction as the network call.}
  \trouble{A gRPC call hangs and the screen never updates}{No deadline was set. Attach one, and handle the deadline-exceeded status instead of swallowing it.}
  \trouble{BLE scan returns no devices}{A runtime permission was never requested: \code{BLUETOOTH\_SCAN} on Android, the usage string on iOS.}
\end{frame}

\closingframe{Push before time runs out.}{Branch lab-7 · one pull request · nothing after the deadline}

\end{document}
